\section{Phase 3: Controlled encoding study and cross-graph check}
\label{sec:phase3-study}
This appendix reports the separate Phase 3A--3F encoding investigation.
It is not pooled with the original 39-graph SAT--ILP repetitions, the
69-graph six-method benchmark, or the later proof-closure counts. Phase 3E
uses the canonical \texttt{\_02} run, and Phase 3F uses the corrected
\texttt{\_01} run;
failed sibling attempts remain excluded.

\subsection{Interruption semantics and construction audits}
The Phase 3A diagnostic pilot contains 30 method-task arms, with recorded
outcomes of 9 \texttt{OPT}, 21 \texttt{UNKNOWN}, and zero \texttt{UNSAT};
27 SAT attempts completed, and five nominal soft-threshold overruns were
recorded. The implementation audit distinguishes internal-deadline
\texttt{UNKNOWN} from parent termination. A safe internal unknown with
final counters was demonstrated for a controlled backend implementing a
tested interrupt contract. For the installed PySAT CaDiCaL195 wrapper,
however, \texttt{interrupt()} and \texttt{clear\_interrupt()} are
unsupported stubs. A real CaDiCaL195 \texttt{solve\_limited()} call did
return \texttt{None} under a conflict budget and exposed accumulated
statistics, but this does not demonstrate active wall-clock interruption.
The benchmark's 30-second parent timeout therefore remains a fallback;
parent-killed attempts retain null final counters and are not represented
as completed internal unknowns.

Phase 3B.1 recorded 60 fixed-span CNF builds. The different ER-60
$(3,2)$ formula sizes across methods were explained by the label-domain
span used to construct each model, rather than by a change in the
Reference encoder. Phase 3C examined 20 fixed-span models across ten
tasks: all 4,660 distance-two pair events were unique, with no duplicate
or tautological distance-two clauses observed. For free endpoints, the
number of forbidden ordered label pairs for a gap $g$ is
\[
 P(U,g)=(2g-1)(U+1)-g(g-1),
\]
when $1\leq g\leq U+1$. Thus $P(U,2)=3U+1$, giving 382 clauses per
distance-two pair at $U=127$. These counts describe the audited
construction and do not imply that fewer clauses improve solver speed.

\subsection{Threshold prototype and fixed-span results}
The Phase 3D isolated prototype passed all 24 pair-equivalence cases
and all eight small complete-model cases under independent arithmetic
validation. It was not integrated into the production solver; its
complete-model tests did not include production symmetry clauses. The
fixed-span comparison is therefore a restricted implementation study,
not evidence of a universal encoding replacement.

The canonical Phase 3E \texttt{\_02} dataset uses the ER-60 graph
\texttt{ER\_60\_p0.30\_seed100}, $(h,k)=(3,2)$, and fixed feasible spans
$U=112$ and $U=127$, with three fresh-process operational runs per
encoding/span. All 12 runs returned SAT with valid witnesses. At
$U=112$, Threshold won one and lost two matched solve-time blocks; at
$U=127$, it was slower in the solve slice in all three blocks, but faster
end-to-end in all three because construction and loading dominated the
short solve. The pairwise formula sizes in the main results table show
fewer clauses and literals but more variables for Threshold. These
fixed-span SAT outcomes do not prove the spans optimal. In particular,
$U=111$ remains unclassified.

The Phase 3E.2 IPASIR-UP pilot found only intrusive, incomplete
instrumentation. Propagator callbacks and observed variables can affect
backend preprocessing and execution; the pilot does not provide
observation-free internal branching traces. Its counters are not treated
as canonical performance evidence or used to infer a causal mechanism.
The final Phase 3E \texttt{\_02} measurements use passive solver statistics.

\subsection{Phase 3F cohort and interpretation}
Phase 3F tests eight selected graphs: two each from the Tree, BA, ER
$p=0.15$, and ER $p=0.30$ families, including one 20- and one
40-vertex graph per family. It uses $(h,k)=(3,2)$, one independently
validated feasible span per graph, two encodings, and three operational
repetitions. The 48 scheduled runs produced 42 SAT results with valid
witnesses, six parent timeouts, zero UNSAT classifications, and zero
worker errors. All six parent timeouts occurred on
\texttt{ER\_40\_p0.30\_seed100} at $U=60$, for both encodings and all
three repetitions. Their final counters are unavailable; they are
right-censored outcomes, not UNSAT.

Across this selected cohort, Threshold reduced clauses by 39.0--43.0\%
and literals by 54.9--58.0\%, while increasing variables. On the seven
graphs with completed runs, its median end-to-end time was lower on six,
but solve-time direction varied by graph; no completed comparison exists
for the dense graph that timed out. The three launches per configuration
are operational repetitions, not independent samples of graph instances.
The cohort contains only eight graphs and one parameter pair, and its
fixed spans are feasible witnesses rather than optimality proofs.

\paragraph{Closure decision.}
Phase 3F concluded with Decision D: the data are inconclusive for a broad
external-validity or practical-speed claim. Retain the Reference encoder
as the production/default path and Threshold as a secondary exploratory
prototype. Phase 3A--3F are closed for the stated research objectives;
no Phase 3G is recommended. These results do not alter the historical
benchmark, its proof certificates, or the original SAT--ILP conclusions.
